% GENERATED FILE. Edit canonical lessons, then run npm run generate:books.
% canonical-source-hash: fnv1a64:8b742824afb3fe31
% canonical-lessons: TE-C466-caurasta, TE-C466-udyanavanam, TE-C466-paccika, TE-C466-vidhidipam, TE-C466-parisaralu

\chapter{Around the Square}
\label{ch:te-around-the-square}
\hlchaptermodality{466}

\begin{chapteropening}
\textbf{By the end of this chapter:} \emph{I can say a town square, a public garden, a lawn, a street lamp and surroundings.}

Complete the last lesson of chapter 466: I can say a town square, a public garden, a lawn, a street lamp and surroundings.
\end{chapteropening}

\section[\texorpdfstring{\te{చౌరస్తా} (caurastā)}{చౌరస్తా (caurastā)}]{\te{చౌరస్తా} (caurastā) — a town square; a crossroads}

\label{lesson:TE-C466-caurasta}



\begin{quote}
Before the new one: say the Telugu for a local person, then the Telugu for a stranger.
\end{quote}

\subsection*{You'll want to know: \te{చౌరస్తా}}
\textbf{\te{చౌరస్తా}} — \emph{caurastā} — \textquotedblleft{}a town square; a crossroads\textquotedblright{}.

\begin{quote}
\te{చౌరస్తాలో} \te{పెద్ద} \te{గడియారం} \te{ఉంది}.
\end{quote}

\begin{quote}
\emph{caurastālō pedda gaḍiyāraṁ undi.}
\end{quote}

\begin{quote}
There is a big clock in the town square.
\end{quote}

Where four roads meet; the busy centre of a small town.

\begin{cousinweb}
\textbf{\te{చౌరస్తా}} (\emph{caurastā}) is usually traced to Persian \emph{chahār}, \textquotedblleft{}four\textquotedblright{}, and \emph{rāsta}, \textquotedblleft{}a road\textquotedblright{}, which came into Telugu through Hindustani: the place where four roads meet.
\end{cousinweb}

\begin{grammarlens}[title={What you've built}]
One more word for getting around.
\end{grammarlens}

\subsection*{Your turn}
\begin{itemize}
\raggedright
  \item \emph{Say it:} \emph{caurastā}
  \item \emph{Say it:} \emph{caurastā}, once more
  \item \emph{Say it:} \emph{caurastālō pedda gaḍiyāraṁ undi.} — There is a big clock in the town square.
  \item \emph{Recall:} say the Telugu for a local person, then the Telugu for a stranger, then say \emph{caurastā} again
\end{itemize}

\begin{tcolorbox}[breakable,colback=teal!4,colframe=teal!35!black,title={Before you move on}]
What does \te{చౌరస్తా} mean? (\textquotedblleft{}A town square; a crossroads\textquotedblright{}.) Say it once more.
\end{tcolorbox}
\section[\texorpdfstring{\te{ఉద్యానవనం} (udyānavanaṁ)}{ఉద్యానవనం (udyānavanaṁ)}]{\te{ఉద్యానవనం} (udyānavanaṁ) — a public garden; a park}

\label{lesson:TE-C466-udyanavanam}



\begin{quote}
Before the new one: say the Telugu for a stranger, then the Telugu for a town square.
\end{quote}

\subsection*{You'll want to know: \te{ఉద్యానవనం}}
\textbf{\te{ఉద్యానవనం}} — \emph{udyānavanaṁ} — \textquotedblleft{}a public garden; a park\textquotedblright{}.

\begin{quote}
\te{సాయంత్రం} \te{ఉద్యానవనంలో} \te{చాలా} \te{మంది} \te{నడుస్తారు}.
\end{quote}

\begin{quote}
\emph{sāyantraṁ udyānavanaṁlō cālā mandi naḍustāru.}
\end{quote}

\begin{quote}
In the evening a lot of people walk in the park.
\end{quote}

\te{ఉద్యానం}, a garden, + \te{వనం}, a grove. \te{పార్కు} is the everyday loan.

\begin{cousinweb}
\textbf{\te{ఉద్యానవనం}} (\emph{udyānavanaṁ}) is a compound of two words borrowed from Sanskrit: \emph{udyāna}, \textquotedblleft{}a park, a garden\textquotedblright{} (CDIAL 2052), and \emph{vana}, \textquotedblleft{}a wood, a grove\textquotedblright{}.
\end{cousinweb}

\begin{grammarlens}[title={What you've built}]
One more word for getting around.
\end{grammarlens}

\subsection*{Your turn}
\begin{itemize}
\raggedright
  \item \emph{Say it:} \emph{udyānavanaṁ}
  \item \emph{Say it:} \emph{udyānavanaṁ}, once more
  \item \emph{Say it:} \emph{sāyantraṁ udyānavanaṁlō cālā mandi naḍustāru.} — In the evening a lot of people walk in the park.
  \item \emph{Recall:} say the Telugu for a stranger, then the Telugu for a town square, then say \emph{udyānavanaṁ} again
\end{itemize}

\begin{tcolorbox}[breakable,colback=teal!4,colframe=teal!35!black,title={Before you move on}]
What does \te{ఉద్యానవనం} mean? (\textquotedblleft{}A public garden; a park\textquotedblright{}.) Say it once more.
\end{tcolorbox}
\section[\texorpdfstring{\te{పచ్చిక} (paccika)}{పచ్చిక (paccika)}]{\te{పచ్చిక} (paccika) — a lawn; grassland}

\label{lesson:TE-C466-paccika}



\begin{quote}
Before the new one: say the Telugu for a town square, then the Telugu for a public garden.
\end{quote}

\subsection*{You'll want to know: \te{పచ్చిక}}
\textbf{\te{పచ్చిక}} — \emph{paccika} — \textquotedblleft{}a lawn; grassland\textquotedblright{}.

\begin{quote}
\te{పిల్లలు} \te{పచ్చిక} \te{మీద} \te{ఆడుకుంటున్నారు}.
\end{quote}

\begin{quote}
\emph{pillalu paccika mīda āḍukuṇṭunnāru.}
\end{quote}

\begin{quote}
The children are playing on the lawn.
\end{quote}

Fresh green grass for sitting or grazing. \te{పచ్చ}, green, is in the word.

\begin{cousinweb}
\textbf{\te{పచ్చిక}} (\emph{paccika}) is native Dravidian: DEDR 3821 lists Telugu \emph{paccika}, \textquotedblleft{}grass\textquotedblright{}, beside \emph{pacca}, \textquotedblleft{}green\textquotedblright{}.
\end{cousinweb}

\begin{grammarlens}[title={What you've built}]
One more word for getting around.
\end{grammarlens}

\subsection*{Your turn}
\begin{itemize}
\raggedright
  \item \emph{Say it:} \emph{paccika}
  \item \emph{Say it:} \emph{paccika}, once more
  \item \emph{Say it:} \emph{pillalu paccika mīda āḍukuṇṭunnāru.} — The children are playing on the lawn.
  \item \emph{Recall:} say the Telugu for a town square, then the Telugu for a public garden, then say \emph{paccika} again
\end{itemize}

\begin{tcolorbox}[breakable,colback=teal!4,colframe=teal!35!black,title={Before you move on}]
What does \te{పచ్చిక} mean? (\textquotedblleft{}A lawn; grassland\textquotedblright{}.) Say it once more.
\end{tcolorbox}
\section[\texorpdfstring{\te{వీధిదీపం} (vīdhidīpaṁ)}{వీధిదీపం (vīdhidīpaṁ)}]{\te{వీధిదీపం} (vīdhidīpaṁ) — a street lamp}

\label{lesson:TE-C466-vidhidipam}



\begin{quote}
Before the new one: say the Telugu for a public garden, then the Telugu for a lawn.
\end{quote}

\subsection*{You'll want to know: \te{వీధిదీపం}}
\textbf{\te{వీధిదీపం}} — \emph{vīdhidīpaṁ} — \textquotedblleft{}a street lamp\textquotedblright{}.

\begin{quote}
\te{మా} \te{వీధిలో} \te{వీధిదీపాలు} \te{వెలగడం} \te{లేదు}.
\end{quote}

\begin{quote}
\emph{mā vīdhilō vīdhidīpālu velagaḍaṁ lēdu.}
\end{quote}

\begin{quote}
The street lamps in our street are not working.
\end{quote}

Plural: \te{వీధిదీపాలు}, \emph{vīdhidīpālu}.

\te{వీధి}, a street, + \te{దీపం}, a lamp. Plural \te{వీధిదీపాలు}.

\begin{cousinweb}
\textbf{\te{వీధిదీపం}} (\emph{vīdhidīpaṁ}) is a compound of two words borrowed from Sanskrit: \emph{vīthi}, \textquotedblleft{}a road, a street\textquotedblright{} (CDIAL 12050), and \emph{dīpa}, \textquotedblleft{}a lamp\textquotedblright{} (CDIAL 6348).
\end{cousinweb}

\begin{grammarlens}[title={What you've built}]
One more word for getting around.
\end{grammarlens}

\subsection*{Your turn}
\begin{itemize}
\raggedright
  \item \emph{Say it:} \emph{vīdhidīpaṁ}
  \item \emph{Say it:} \emph{vīdhidīpaṁ}, once more
  \item \emph{Say it:} \emph{mā vīdhilō vīdhidīpālu velagaḍaṁ lēdu.} — The street lamps in our street are not working.
  \item \emph{Recall:} say the Telugu for a public garden, then the Telugu for a lawn, then say \emph{vīdhidīpaṁ} again
\end{itemize}

\begin{tcolorbox}[breakable,colback=teal!4,colframe=teal!35!black,title={Before you move on}]
What does \te{వీధిదీపం} mean? (\textquotedblleft{}A street lamp\textquotedblright{}.) Say it once more.
\end{tcolorbox}
\section[\texorpdfstring{\te{పరిసరాలు} (parisarālu)}{పరిసరాలు (parisarālu)}]{\te{పరిసరాలు} (parisarālu) — surroundings}

\label{lesson:TE-C466-parisaralu}



\begin{quote}
Before the new one: say the Telugu for a lawn, then the Telugu for a street lamp.
\end{quote}

\subsection*{You'll want to know: \te{పరిసరాలు}}
\textbf{\te{పరిసరాలు}} — \emph{parisarālu} — \textquotedblleft{}surroundings\textquotedblright{}.

\begin{quote}
\te{మన} \te{పరిసరాలను} \te{శుభ్రంగా} \te{ఉంచాలి}.
\end{quote}

\begin{quote}
\emph{mana parisarālanu śubhraṅgā uñcāli.}
\end{quote}

\begin{quote}
We should keep our surroundings clean.
\end{quote}

Mostly plural. The area around a house, a school or a town.

\begin{cousinweb}
\textbf{\te{పరిసరాలు}} (\emph{parisarālu}) is borrowed from Sanskrit \emph{parisara}, \textquotedblleft{}the surroundings, the neighbourhood\textquotedblright{} (Monier-Williams), with the Telugu plural -\te{లు}.
\end{cousinweb}

\begin{grammarlens}[title={What you've built}]
One more word for getting around.
\end{grammarlens}

\subsection*{Your turn}
\begin{itemize}
\raggedright
  \item \emph{Say it:} \emph{parisarālu}
  \item \emph{Say it:} \emph{parisarālu}, once more
  \item \emph{Say it:} \emph{mana parisarālanu śubhraṅgā uñcāli.} — We should keep our surroundings clean.
  \item \emph{Recall:} say the Telugu for a lawn, then the Telugu for a street lamp, then say \emph{parisarālu} again
\end{itemize}

\begin{tcolorbox}[breakable,colback=teal!4,colframe=teal!35!black,title={Before you move on}]
What does \te{పరిసరాలు} mean? (\textquotedblleft{}Surroundings\textquotedblright{}.) Say it once more.
\end{tcolorbox}
